\chapter{Conclusioni}

\section{Sintesi del caso studio}

Il presente lavoro di tesi ha avuto come obiettivo principale la progettazione, implementazione e valutazione di un sistema di \emph{Anomaly Detection} basato sull’algoritmo Isolation Forest, applicato all’analisi di log HTTP provenienti da un contesto reale.

Il punto di partenza dell’elaborato risiede nella constatazione che i sistemi di difesa tradizionali, fondati su firme statiche e regole predefinite, mostrano limiti significativi nel fronteggiare la crescente dinamicità delle minacce informatiche e l’elevata variabilità del traffico web moderno. In tale scenario, approcci basati su tecniche di Machine Learning risultano particolarmente promettenti, in quanto permettono di individuare comportamenti anomali senza la necessità di una conoscenza preventiva delle specifiche tipologie di attacco.

L’attività svolta si è quindi concentrata sulla costruzione di un dataset derivato da log HTTP, opportunamente trasformato attraverso un processo di \emph{feature engineering} volto a estrarre metriche significative dal punto di vista della sicurezza applicativa. Su tali dati è stato applicato un modello non supervisionato, con l’obiettivo di verificare la capacità dell’algoritmo di identificare automaticamente richieste anomale o potenzialmente malevole.

I risultati ottenuti evidenziano come l’approccio adottato sia efficace nell’individuare comportamenti anomali all’interno del traffico analizzato ma che risulti non direzionato all'ambito cybersecurity. In particolare, il modello dimostra una buona sensibilità nel rilevare pattern atipici, confermando la validità dell’ipotesi iniziale. Tuttavia, come prevedibile nei sistemi di anomaly detection non supervisionati, emerge la presenza di falsi positivi, spesso associati a traffico lecito ma caratterizzato da elevata complessità o intensità, oscurando del tutto quelli che possono essere attacchi reali.

Tale comportamento evidenzia un compromesso intrinseco tra sensibilità e precisione: se da un lato il modello è in grado di intercettare un ampio spettro di anomalie, dall’altro richiede ulteriori meccanismi di filtraggio o validazione per ridurre il numero di segnalazioni non rilevanti a beneficio di quelli realmente rilevanti per l'utente finale. Questo aspetto rappresenta una caratteristica comune nei modelli di Machine Learning applicati alla cybersecurity, in cui il raggiungimento simultaneo di elevati livelli di precisione e accuratezza risulta particolarmente complesso, soprattutto in assenza di dati etichettati.

\section{Considerazioni finali}
